\section{Spatial derivatives and exact harmonic contact terms}\label{tsc:contacts}
The identity $\partial_x\calB_s=\calB_{s-1}$ is distributional. Passing
from it to a pointwise polygamma product requires a local correction. This
section gives that correction explicitly and does not alter the correct
pointwise derivative tower already present in the repository.

\subsection{Finite reductions at negative integral spectral orders}
For $r\ge1$, $\calB_{-r}=\delta_0^{(r)}$. Applying this to the master
identity gives a useful family of exact spectral resonances.
\begin{proposition}\label{tsc:negative-resonance}
For $r\ge1$ and any complex $s,t$,
\begin{align}\label{tsc:negative-resonance-eq}
 \calT(-r,s,t;\mathbf a)
 ={}&(2\pi)^{s+t-r}(-1)^r
   \sum_{j=0}^r\binom rj
    \calB_{s-j}(a_1-a_2)\calB_{t-r+j}(a_1-a_3).
\end{align}
Here the functions on the right are evaluated away from their seam using
\eqref{tsc:Bhurwitz}. At order zero, the missing constant mode changes
the formula to
\begin{align}\label{tsc:zero-resonance}
 \calT(0,s,t;\mathbf a)={}&(2\pi)^{s+t}
    \calB_s(a_1-a_2)\calB_t(a_1-a_3)\notag\\
 &-e^{-\pi\ii(s-t)/2}\Li_{s+t}(e(a_3-a_2))\notag\\
 &-e^{\pi\ii(s-t)/2}\Li_{s+t}(e(a_2-a_3)).
\end{align}
\end{proposition}
\begin{proof}
Pair $\delta_{a_1}^{(r)}$ with the product of the other two, smooth,
factors. The sign $(-1)^r$ is the defining sign of a delta derivative,
and the product rule with $\partial_x^j\calB_s=\calB_{s-j}$ gives
\eqref{tsc:negative-resonance-eq}. For zero order use
$\calB_0=\delta_0-1$, not $\delta_0$. The remaining bilinear mean is
computed by its two opposite-frequency sectors and equals the two
polylogarithm terms after multiplication by $(2\pi)^{s+t}$.
\end{proof}
The formulas reduce the complete six-sector expression on these spectral
hyperplanes. They do not determine its transverse derivatives, which are
precisely among the data needed for general Stieltjes triple products.

\subsection{Elementary harmonic coefficients}
Define $e_j^{(r)}$ by
\begin{equation}\label{tsc:elementary}
 \prod_{k=1}^r(1+v/k)=\sum_{j=0}^r e_j^{(r)}v^j,
 \qquad e_j^{(r)}=0\quad(j>r).
\end{equation}
Thus $e_1^{(r)}=H_r$,
$e_2^{(r)}=(H_r^2-H_r^{(2)})/2$, and
$e_3^{(r)}=(H_r^3-3H_rH_r^{(2)}+2H_r^{(3)})/6$.
Unlike the complete coefficients in the primitive ladder, these have finite
support in $j$.

Let $D_{m,r}$ denote the scale-one, one-sided Hadamard extension of the
\emph{pointwise} function $\gamma_m^{(r)}(x)$ on $0<x<1$. At the right-hand
seam, expand the test function through degree $r$, integrate all singular
power-log terms, and delete the divergent terms in the cutoff. This
specifies the extension without arbitrary delta terms. Put $D_{m,0}=C_m$.

\begin{theorem}[Derivative--extension conversion]\label{tsc:contact-single}
For $r\ge0$,
\begin{equation}\label{tsc:single-contact}
 \partial_x^r C_m=D_{m,r}+c_{m,r}\delta_0^{(r)},
 \qquad c_{m,r}=(-1)^{m+1}m!e_{m+1}^{(r)}.
\end{equation}
In particular $c_{m,0}=0$, $c_{0,r}=-H_r$, and $c_{m,r}=0$ for $m\ge r$.
\end{theorem}
\begin{proof}
In a local right-hand coordinate, the meromorphic distribution
$x_+^{u-r-1}$ has residue
$(-1)^r\delta_0^{(r)}/r!$ at $u=0$. Its regular Laurent coefficients
are the scale-one finite parts of $x^{-r-1}\log^j x/j!$.
Differentiating the singular term $x^{u-1}$ in
$\zeta(1-u,x)$ gives
\[
 \partial_x^r x^{u-1}=(-1)^r r!P_r(u)x^{u-r-1},
 \qquad P_r(u)=\prod_{k=1}^r(1-u/k).
\]
Hence the residue contribution after multiplication is
$P_r(u)\delta_0^{(r)}/u$. The regular nonresidue terms give exactly the
raw extensions $D_{m,r}$. Comparing with the $r$th derivative of
\eqref{tsc:Zlaurent}, the additional coefficient is
$m![u^{m+1}]P_r(u)=(-1)^{m+1}m!e_{m+1}^{(r)}$.
The remaining smooth local terms introduce no contact distribution.
For $r=0$ the assertion is immediate. This proves the formula.
\end{proof}

\begin{theorem}[Contact-complete derivative triple]\label{tsc:contact-triple}
Let $r_i,m_i\ge0$, $R=r_1+r_2+r_3$, and retain distinct shifts. Then
\begin{align}\label{tsc:contact-triple-eq}
 &\FP\int_0^1\prod_{i=1}^3
            \gamma_{m_i}^{(r_i)}(\{x-a_i\})\dd x\notag\\
 &\quad=(2\pi)^R\left(\prod_{i=1}^3m_i!\right)
    [\mathbf u^{\mathbf m+\mathbf1}]
        \left(\prod_{i=1}^3E(u_i)\right)
                       \calT(\mathbf u-\mathbf r;\mathbf a)\notag\\
 &\qquad\quad-
    \sum_{i=1}^3(-1)^{r_i}c_{m_i,r_i}
       \left.\partial_x^{r_i}
              \prod_{j\ne i}\gamma_{m_j}^{(r_j)}(\{x-a_j\})
       \right|_{x=a_i}.
\end{align}
The local derivatives in the last line are ordinary derivatives of smooth
functions at $a_i$. In particular, setting all $m_i=0$ gives identities for
arbitrary products of three separated polygamma functions.
\end{theorem}
\begin{proof}
By \eqref{tsc:single-contact}, each raw factor is
$D_{m_i,r_i,a_i}=\partial_x^{r_i}C_{m_i,a_i}
                 -c_{m_i,r_i}\delta_{a_i}^{(r_i)}$.
In the product expansion, every term containing two delta-supported
factors is zero because their supports are disjoint. A term containing
one delta derivative pairs with the smooth remaining product, producing
$(-1)^{r_i}$ times its ordinary $r_i$th derivative at $a_i$. These are
exactly the terms in the final line of \eqref{tsc:contact-triple-eq}.

The product of the three periodic derivatives is
$(-1)^R\partial_{a_1}^{r_1}\partial_{a_2}^{r_2}\partial_{a_3}^{r_3}
\calJ_{\mathbf m}$. Apply \eqref{tsc:lowering} to
\eqref{tsc:all-index}. The factor $(-1)^R$ cancels the sign of
$(-2\pi)^R$, giving the first line of the right-hand side.
\end{proof}

The correction is not optional. For instance,
$\partial_x C_0=D_{0,1}-\delta_0'$. Omitting it already changes a triple
with one spatial derivative by the derivative of the product of the two
other endpoint values. The theorem explains the correction without
claiming that the manuscript's pointwise formulas are wrong.

For explicit simplification of endpoint derivatives, the classical
Pochhammer identity yields, for $r\ge1$,
\begin{equation}\label{tsc:pointwise-tower}
 \gamma_n^{(r)}(x)=(-1)^{r+n}r!n!
      \sum_{k=0}^{\min(r,n)}
       e_k^{(r)}\frac{\zeta^{[n-k]}(r+1,x)}{(n-k)!}.
\end{equation}
Indeed, expand
$\partial_x^r\zeta(1+v,x)=(-1)^r(1+v)_r\zeta(1+r+v,x)$ at zero.
This derivation also fixes every factorial and sign in
\eqref{tsc:pointwise-tower}.
